\documentclass[11pt]{article}
\usepackage{fullpage,amssymb,graphicx,hyperref,xcolor}

\newcommand{\fig}[2]{\begin{figure}[ht]\begin{center}%
  \includegraphics[scale=#2]{#1}\end{center}%
  \end{figure}}
\newcommand{\ceil}[1]{\left\lceil #1\right\rceil}

\newcommand{\Ex}{\mathbb{E}}
\newcommand{\Var}{\textrm{Var}}

\begin{document}

\begin{center}\Large\bf 
CS 473 (Spring 2025)\\
Homework 6 (due Mar 27 Thu 10am)
\end{center}

\medskip\noindent
{\bf Instructions:} As in previous homeworks.

\begin{description}
\bigskip
\item[Problem 6.1:]
The original version of the smallest enclosing circle problem
is very sensitive to so-called ``outliers'': if a single
point is corrupted, the optimal circle could change drastically.
This motivates the following extension of the problem:
\begin{quote}
Given a set $S$ of $n$ points in 2D and an integer $k$, find the
smallest circle that encloses at least $n-k$ points (i.e.,
we allow at most $k$ points to be outside the circle).
\end{quote}
In this question, you will design an efficient Monte Carlo
algorithm to solve this new problem.  You may assume that there are no
degeneracies, i.e., no 4 points lie on a common circle.

\begin{enumerate}
\item[(a)] (10 pts) Describe a naive algorithm that runs in $O(n^4)$ time.

\smallskip
\item[(b)] (35 pts) Let $C^*$ denote the optimal circle, let $B^*$
denote the set of the (at most) 3 points on the boundary of $C^*$,
 and 
let $Q^*$ denote the set of at most $k$ points outside the
circle.  
Form a random subset $R\subseteq S$ as follows:
for each point $p\in S$, put $p$ in $R$ independently with probability $1/k$.
(Consequently, $\Ex[|R|]=n/k$.)  

Prove that the probability
that $B^*\subseteq R\subseteq S-Q^*$ is at least $\Omega(1/k^3)$.

\smallskip
\item[(c)] (40 pts) Describe a Monte Carlo algorithm that runs in
$O(n)$ worst-case time and is correct with probability 
$\Omega(1/k^3)$.
(Yes, this probability converges to 0, i.e., your
algorithm is supposed to wrong most of the time!)

[Hint: you may use the smallest enclosing circle algorithm
from class as a subroutine; 
obviously, (b) is meant to help\ldots]

\smallskip
\item[(d)] (15 pts) Now, describe an efficient Monte Carlo algorithm that has error probability at most $0.01$.  Analyze the worst-case running time as a function of $n$ and $k$.  When $k$ is small, it should be much faster than the algorithm in (a).
\end{enumerate}



\bigskip
\item[Problem 6.2:]
\newcommand{\AREA}{\mathop{\rm area}}
We are given $n$ triangles $T_1,\ldots,T_n$ in 2D\@.
We want to compute the area of their union, i.e., $A=\AREA(T_1\cup\cdots\cup T_n)$.
(Note that the triangles may overlap, so $A$ may not necessarily be equal to
$\AREA(T_1)+\cdots+\AREA(T_n)$.)

\fig{hw6fig1.jpg}{0.9}

In this problem, we consider a randomized Monte Carlo algorithm to compute a good approximation of $A$.
Let $r$ be a parameter.  Here is the pseudocode:
\begin{quote}
\begin{tabbing}
1. \= for \=\kill
1.\> for $i=1$ to $n$ do\\
2.\>\> independently pick $r$ random points $p_{i,1},\ldots,p_{i,r}$ from $T_i$\\
3.\> return $\displaystyle X \:=\: \sum_{i=1}^n |\{p_{i,1},\ldots,p_{i,r}\}\setminus (T_1\cup\cdots\cup T_{i-1})|\cdot \frac{\AREA(T_i)}{r}$
\end{tabbing}
\end{quote}

\begin{enumerate}
\item[(a)] (30 pts) Prove that $\Ex[X] = A$.

[Hint: what is $\sum_{i=1}^n \AREA(T_i\setminus (T_1\cup\cdots\cup T_{i-1}))$?]

\smallskip
\item[(b)] (30 pts) Prove that $\Var[X] \le {\color{red}A^2/r}$.  ($\Var[X]$ denotes the variance of $X$.)

\smallskip
\item[(c)] (40 pts) By choosing $r$ appropriately, show that the algorithm runs in {\color{red}$O(n^2)$} time and
outputs a value $X$ between $0.99A$ and $1.01A$ with probability at least 0.99.  

[Note: you may assume that we can generate a random point in a triangle in $O(1)$ time,
we can compute the area of a triangle in $O(1)$ time, and we can test whether a point is inside a triangle in $O(1)$ time.  The above algorithm actually works for other types of objects besides triangles, and in dimension higher than 2.]
\end{enumerate}


\bigskip
\item[Problem 6.3:]
We are given a (unweighted) bipartite graph $G$ with $n$ vertices and $m$ edges.  Suppose we have already computed a maximum matching $M$ in $G$.  Now suppose we delete
a vertex $v$ from $G$ and all its incident edges.  Let $G'$
be the new graph (with $n-1$ vertices).
Describe how to efficiently compute a new maximum matching $M'$
in $G'$.  (Your algorithm should be faster than re-running a matching
algorithm from scratch.  Remember to prove correctness of your algorithm.)



\end{description}
\end{document}
